\lecture{4.3}{G-Computation Using Regression Models}
\label{lec:4-3}

\subsection{Prediction-Based Implementation}

Empirical g-computation by hand is transparent, but it becomes cumbersome when $W$
has many components. Regression provides a systematic way to estimate
$\Qbar(a,w)=\E(Y\given A=a,W=w)$ and then evaluate the fitted regression under
counterfactual treatment values.

The generic regression-based g-computation estimator is
\[
\psi_n(a)=\frac{1}{n}\sum_{i=1}^n \Qbar_n(a,W_i),
\]
where $\Qbar_n$ is a fitted outcome regression. The estimated ATE is
\[
\widehat{\ATE}_{\mathrm{gcomp}}
=
\frac{1}{n}\sum_{i=1}^n
\left\{\Qbar_n(1,W_i)-\Qbar_n(0,W_i)\right\}.
\]

This estimator has two conceptually distinct stages. First, estimate conditional
outcome means from the observed data. Second, average predicted outcomes under
treatment levels that may differ from the treatment actually observed for a given
individual.

\subsection{Using \texttt{predict()} with Counterfactual Treatment Values}

In software, the most stable way to implement g-computation is to create two
modified copies of the covariate data:
\[
\text{data}^{(1)}=(W_i,A_i=1)_{i=1}^n,\qquad
\text{data}^{(0)}=(W_i,A_i=0)_{i=1}^n.
\]
The fitted outcome model is then used to predict the outcome for every row of each
modified dataset.

\begin{implementationbox}
Generic R pattern for g-computation:
\begin{lstlisting}[language=R]
# fit an outcome regression
fit_Q <- glm(Y ~ A + W, data = dat, family = gaussian())

# create treatment-specific prediction datasets
dat1 <- dat
dat1$A <- 1
dat0 <- dat
dat0$A <- 0

# predict under each treatment level
Q1 <- predict(fit_Q, newdata = dat1, type = "response")
Q0 <- predict(fit_Q, newdata = dat0, type = "response")

# estimate counterfactual means and ATE
psi1 <- mean(Q1)
psi0 <- mean(Q0)
ate  <- psi1 - psi0
\end{lstlisting}
\end{implementationbox}

The argument \texttt{newdata} is crucial. It asks the fitted regression to evaluate
the conditional mean at covariate-treatment combinations chosen by the analyst, not
only at the treatment values actually observed.

\subsection{Saturated GLMs}

For discrete covariates, a saturated generalized linear model can reproduce the
empirical subgroup means. With binary $A$ and binary $W$, a saturated linear model is
\[
\E(Y\given A=a,W=w)
=
\beta_0+\beta_1a+\beta_2w+\beta_3aw.
\]
If a Gaussian working model with identity link is fit by least squares, the predictions
at the four $(a,w)$ combinations equal the four observed subgroup sample means.

For a binary outcome, a saturated logistic model can be written as
\[
\logit\{\Prob(Y=1\given A=a,W=w)\}
=
\beta_0+\beta_1a+\beta_2w+\beta_3aw.
\]
Its fitted probabilities can match the empirical risks in each subgroup, provided the
data do not have complete separation and each subgroup contains observations. The
link function changes the parameterization, but saturation still means that no
subgroup restriction is imposed beyond the chosen covariate coding.

\subsection{Main-Terms GLMs}

In practice, analysts often use smoother, lower-dimensional regressions. A main-terms
linear outcome model is
\[
\Qbar_n(a,w)=\hat\beta_0+\hat\beta_1 a+\hat\beta_2^\top w.
\]
The g-computation estimator is then
\[
\psi_n(a)
=
\frac{1}{n}\sum_{i=1}^n
\left(\hat\beta_0+\hat\beta_1 a+\hat\beta_2^\top W_i\right)
=
\hat\beta_0+\hat\beta_1 a+\hat\beta_2^\top \bar W_n,
\]
where $\bar W_n=n^{-1}\sum_i W_i$. Therefore,
\[
\widehat{\ATE}_{\mathrm{gcomp}}
=
\psi_n(1)-\psi_n(0)
=
\hat\beta_1.
\]
This equality is a special property of the identity-link linear main-terms model. It
does not hold for most other models.

If the model includes a treatment-covariate interaction,
\[
\Qbar_n(a,w)=\hat\beta_0+\hat\beta_1 a+\hat\beta_2^\top w+\hat\beta_3^\top(aw),
\]
then
\[
\widehat{\ATE}_{\mathrm{gcomp}}
=
\frac{1}{n}\sum_{i=1}^n
\left(\hat\beta_1+\hat\beta_3^\top W_i\right)
=
\hat\beta_1+\hat\beta_3^\top \bar W_n.
\]
The coefficient $\hat\beta_1$ is the estimated treatment contrast when $W=0$, not the
population-average treatment effect unless $\bar W_n=0$ or $\hat\beta_3=0$.

\subsection{Logistic Regression G-Computation}

Suppose $Y$ is binary and we fit a main-terms logistic regression
\[
\Qbar_n(a,w)
=
\expit(\hat\gamma_0+\hat\gamma_1a+\hat\gamma_2^\top w).
\]
The g-computation estimator is
\[
\psi_n(1)
=
\frac{1}{n}\sum_{i=1}^n
\expit(\hat\gamma_0+\hat\gamma_1+\hat\gamma_2^\top W_i),
\]
\[
\psi_n(0)
=
\frac{1}{n}\sum_{i=1}^n
\expit(\hat\gamma_0+\hat\gamma_2^\top W_i),
\]
and
\[
\widehat{\ATE}_{\mathrm{gcomp}}
=
\frac{1}{n}\sum_{i=1}^n
\left[
\expit(\hat\gamma_0+\hat\gamma_1+\hat\gamma_2^\top W_i)
-
\expit(\hat\gamma_0+\hat\gamma_2^\top W_i)
\right].
\]

\begin{implementationbox}
Logistic regression g-computation in R:
\begin{lstlisting}[language=R]
fit_Q <- glm(Y ~ A + W, data = dat, family = binomial())

dat1 <- dat
dat1$A <- 1
dat0 <- dat
dat0$A <- 0

Q1 <- predict(fit_Q, newdata = dat1, type = "response")
Q0 <- predict(fit_Q, newdata = dat0, type = "response")

psi1 <- mean(Q1)
psi0 <- mean(Q0)
ate  <- psi1 - psi0
\end{lstlisting}
\end{implementationbox}

\subsection{Why Logistic Coefficients Are Not ATEs}

The coefficient $\hat\gamma_1$ in a logistic regression is a conditional log odds ratio
parameter. It compares the log odds of the outcome between treatment groups at a
fixed value of $W$, assuming the model is correctly specified on the log odds scale.
The ATE is a marginal risk difference:
\[
\E\{\Yone\}-\E\{\Yzero\}.
\]
These are different objects.

There are three reasons the logistic treatment coefficient is not generally an ATE.
First, it is on the log odds scale, whereas the ATE is a mean difference or risk
difference. Second, it is conditional on $W$, whereas the ATE averages over the
population distribution of $W$. Third, the expit function is nonlinear, so averaging
after applying the inverse-logit is not the same as applying the inverse-logit after
averaging:
\[
\frac{1}{n}\sum_{i=1}^n \expit(\eta_i)
\neq
\expit\left(\frac{1}{n}\sum_{i=1}^n \eta_i\right)
\]
in general.

\begin{keyidea}
For regression-based g-computation, the causal estimate is the average contrast of
predicted counterfactual outcomes. Regression coefficients are intermediate
nuisance quantities. They are not automatically causal effects.
\end{keyidea}
